% Unit 131 terjemahan Bahasa Indonesia dari chapter9.tex baris 1626--1722.
% Sumber: Methods of Algebra, Volume 2, commit
% 9a5803ff77dd3257484cb177f851a73770a59dd3 (CC BY 4.0).
% stable-unit-id: o014.aljabr2.chapter9.tannakian-a
% Cakupan: Bagian 9.7 dari definisi kategori Tannakian hingga contoh gradasi.

% segment-id: o014.li2.chapter9.tannakian-a.h001
\section{Kategori Tannakian}\label{sec:Tannakian-cat}
\hypertarget{o014.aljabr2.chapter9.tannakian-a}{ }
% segment-id: o014.li2.chapter9.tannakian-a.p001
Teori kategori Tannakian merupakan perumusan ulang menyeluruh oleh Saavedra
Rivano dan Deligne \cite{Sa72,Del90} atas karya Tadao Tannaka \cite{Ta38}
dan M.\ Krein \cite{Kr49} mengenai representasi grup kompak, dengan gagasan
yang berasal dari Grothendieck. Gagasan terkait kemudian diperluas secara
mendalam dalam karya Lurie dan Bhatt tentang geometri aljabar turunan. Kita
akan mengikuti pendekatan Deligne secara garis besar.

% segment-id: o014.li2.chapter9.tannakian-a.p002
Kecuali dinyatakan lain, sepanjang bagian ini kita mempertahankan asumsi
dari \S\ref{sec:Tannaka-coend}, dengan memilih medan $\Bbbk$.

% segment-id: o014.li2.chapter9.tannakian-a.d001
\begin{definition}\label{def:Tannakian-cat}
	\index{kategori Tannakian@kategori Tannakian (Tannakian category)}
	\index{funktor serat@funktor serat (fiber functor)}
	Misalkan $\mathcal{T}$ kategori abelian kecil secara esensial
	(Definisi~\ref{def:essentially-small-cat}) yang juga memiliki struktur
	monoidal simetris sedemikian rupa sehingga $\otimes$ bersifat
	$\Bbbk$-linear dalam setiap variabel. Kategori $\mathcal{T}$ disebut
	\emph{kategori Tannakian di atas medan $\Bbbk$} jika memenuhi
	sifat-sifat berikut.
	\begin{itemize}
		\item Jika $\munit$ unit $\mathcal{T}$, maka
		$\End_{\mathcal{T}}(\munit)=\Bbbk$.
		\item Sebagai kategori monoidal simetris, $\mathcal{T}$ kaku
		(Definisi~\ref{def:rigid-cat-braided}).
		\item Terdapat aljabar-$\Bbbk$ komutatif $B$, yang menurut konvensi tak
		nol, dan funktor monoidal eksak kanan yang kompatibel dengan struktur
		kepang
		% segment-id: o014.li2.chapter9.tannakian-a.q001
		\[ \omega: \mathcal{T} \to \cated{Mod}B. \]
		Funktor semacam ini disebut \emph{funktor serat} bagi $\mathcal{T}$ di atas
		$B$.\footnote{Istilah ini berasal dari aliran Grothendieck dan berakar
		pada suatu analogi dengan teori ruang penutup.}
	\end{itemize}

	Funktor serat di atas $\Bbbk$ disebut \emph{netral}. Jika kategori
	Tannakian $\mathcal{T}$ memiliki funktor serat netral, maka
	$\mathcal{T}$ disebut \emph{kategori Tannakian netral}.
\end{definition}

% segment-id: o014.li2.chapter9.tannakian-a.p003
Menurut definisi morfisme dual dalam Definisi~\ref{def:dual-morphism},
funktor dual $(\mathord\cdot)^*$ dari kategori Tannakian pasti
$\Bbbk$-linear.

% segment-id: o014.li2.chapter9.tannakian-a.p004
Definisi di atas cukup lentur. Jika $\omega$ funktor serat di atas $B$
dan $B\to B'$ homomorfisme aljabar-$\Bbbk$ komutatif, maka
$\omega\dotimes{B}B'$ adalah funktor serat di atas $B'$. Aksioma-aksioma
tersebut memiliki sejumlah konsekuensi tak trivial, dengan dualitas
memainkan peran kunci.

% segment-id: o014.li2.chapter9.tannakian-a.l001
\begin{lemma}\label{prop:Tannaka-generalities}
	Misalkan $\mathcal{T}$ kategori Tannakian.
	\begin{enumerate}[(i)]
		\item Untuk setiap aljabar-$\Bbbk$ komutatif $B$, setiap funktor serat
		$\omega$ di atas $B$ selalu mengambil nilai dalam
		$\catesubd{Mod}{\mathrm{pfg}}B$ dan secara otomatis setia eksak.
		\item Terdapat perluasan medan $\Bbbk'|\Bbbk$ dan funktor serat di
		atas $\Bbbk'$.
		\item Jika $\omega$ funktor serat di atas $B$, pemetaan kanonik
		% segment-id: o014.li2.chapter9.tannakian-a.q002
		\[ \Hom_{\mathcal{T}}(X,Y)\dotimes{\Bbbk}B\to
		\Hom_{\cated{Mod}B}(\omega(X),\omega(Y)) \]
		bersifat injektif.
		\item Kategori $\mathcal{T}$ hingga secara lokal
		(Definisi~\ref{def:locally-finite}).
	\end{enumerate}
\end{lemma}
\begin{proof}
	Untuk (i), karena $\omega$ funktor monoidal dan $\mathcal{T}$ kaku,
	Proposisi~\ref{prop:dualizable-module} menunjukkan bahwa $\omega$
	mengambil nilai dalam $\catesubd{Mod}{\mathrm{pfg}}B$. Berikutnya kita
	buktikan bahwa $\omega$ eksak kiri. Tinjau sembarang barisan eksak dalam
	$\mathcal{T}$
	% segment-id: o014.li2.chapter9.tannakian-a.q003
	\[ 0 \to X \to Y \to Z. \]
	Dualnya juga eksak, sehingga penerapan $\omega$ memberikan barisan eksak
	dalam $\cated{Mod}B$
	% segment-id: o014.li2.chapter9.tannakian-a.q004
	\[ \omega(Z)^\vee \to \omega(Y)^\vee \to \omega(X)^\vee \to 0. \]
	Ambil kembali
	$(\mathord\cdot)^\vee=\Hom_{\cated{Mod}B}(\mathord\cdot,B)$ pada
	barisan ini. Dengan Proposisi~\ref{prop:Hom-left-exact} dan
	$\omega(\mathord\cdot)^{\vee\vee}\simeq\omega(\mathord\cdot)$,
	diperoleh barisan eksak
	% segment-id: o014.li2.chapter9.tannakian-a.q005
	\[ 0 \to \omega(X) \to \omega(Y) \to \omega(Z). \]

	Untuk kesetiaan, definisi dual menunjukkan bahwa $X\neq0$ jika dan
	hanya jika $\mathrm{ev}_X:X\otimes X^*\to\munit$ tak nol. Sementara itu,
	Proposisi~\ref{prop:rigid-unit-simple} dan
	$\End_{\mathcal{T}}(\munit)=\Bbbk$ menunjukkan bahwa ini setara dengan
	$\mathrm{ev}_X$ surjektif. Karena $\omega$ monoidal sekaligus eksak,
	$X\neq0$ mengakibatkan epimorfisme
	$\omega(X)\dotimes{B}\omega(X)^\vee\twoheadrightarrow B$, sehingga
	$\omega(X)\neq0$.

	Untuk (ii), pilih funktor serat $\omega$ di atas $B$ dan ideal
	maksimal $\mathfrak m$ dari $B$, lalu definisikan
	$\Bbbk':=B/\mathfrak m$. Maka $\omega\dotimes{B}\Bbbk'$ memberikan
	funktor serat di atas $\Bbbk'$.

	Untuk (iii), ambil unsur-unsur yang bebas linear
	$f_1,\ldots,f_n\in\Hom_{\mathcal{T}}(X,Y)$. Unsur-unsur tersebut
	mendefinisikan morfisme
	$\Phi:\munit^{\oplus n}\to\iHom(X,Y)$, dengan notasi dari
	Korolari~\ref{prop:rigid-Hom-internal}. Kita tunjukkan bahwa $\Phi$
	monik. Memang, $\munit$ objek sederhana
	(Proposisi~\ref{prop:rigid-unit-simple}), sehingga $\Ker(\Phi)$ objek
	berpanjang hingga. Jika $\Ker(\Phi)\neq0$, semua faktor komposisinya
	ialah $\munit$; lihat \S\ref{sec:semisimple}. Khususnya, terdapat
	monomorfisme
	$\munit\hookrightarrow\Ker(\Phi)\hookrightarrow\munit^{\oplus n}$.
	Karena $\End_{\mathcal{T}}(\munit)=\Bbbk$, komposisinya dengan $\Phi$
	memberikan morfisme $\munit\to\iHom(X,Y)$ yang di satu sisi nol, tetapi
	di sisi lain bersesuaian dengan
	$\sum_{i=1}^na_if_i\in\Hom_{\mathcal{T}}(X,Y)$, dengan
	$a_1,\ldots,a_n\in\Bbbk$ tidak semuanya nol. Ini bertentangan dengan
	kebebasan linear.

	Menerapkan funktor monoidal eksak $\omega$ pada $\Phi$ memberikan
	monomorfisme modul-$B$
	$B^{\oplus n}\hookrightarrow
	\Hom_{\cated{Mod}B}(\omega(X),\omega(Y))$. Ini tepat merupakan restriksi
	morfisme kanonik yang dibahas ke
	$\bigoplus_{i=1}^n\Bbbk f_i\dotimes{\Bbbk}B$.

	Untuk (iv), mula-mula gunakan (ii) untuk memperoleh funktor serat
	$\omega$ di atas suatu perluasan medan $\Bbbk'$. Sifat setia eksak dari
	(i) sudah mengakibatkan setiap $X\in\Obj(\mathcal{T})$ berpanjang
	hingga. Selanjutnya, (iii) memberikan pembenaman $\Bbbk'$-linear
	% segment-id: o014.li2.chapter9.tannakian-a.q006
	\[ \Hom_{\mathcal{T}}(X,Y)\dotimes{\Bbbk}\Bbbk'
	\hookrightarrow\Hom_{\cate{Vect}_{\mathrm f}(\Bbbk')}
	(\omega(X),\omega(Y)). \]
	Ruas kanan berdimensi hingga, sehingga
	$\Hom_{\mathcal{T}}(X,Y)$ berdimensi hingga sebagai ruang
	vektor-$\Bbbk$.
\end{proof}

% segment-id: o014.li2.chapter9.tannakian-a.p005
Berdasarkan Lema~\ref{prop:Tannaka-generalities}, kerangka teoretis
\S\ref{sec:Tannaka-coend} berlaku untuk setiap kategori Tannakian
$\mathcal{T}$ dan setiap funktor serat
$\omega:\mathcal{T}\to\catesubd{Mod}{\mathrm{pfg}}B$ di atas
aljabar-$\Bbbk$ komutatif $B$. Dari sudut pandang geometri aljabar,
$\catesubd{Mod}{\mathrm{pfg}}B$ bersesuaian dengan kategori berkas vektor
di atas skema afin $\Spec B$. Karena itu, meninjau $B$ umum memang perlu,
tetapi juga memerlukan alat teoretis tambahan. Teorema rekonstruksi berikut
hanya membahas kasus $B=\Bbbk$.

% segment-id: o014.li2.chapter9.tannakian-a.th001
\begin{theorem}\label{prop:Tannaka-neutral}
	\index{teorema rekonstruksi@teorema rekonstruksi}
	Misalkan $\mathcal{T}$ kategori Tannakian netral dan
	$\omega:\mathcal{T}\to\cate{Vect}_{\mathrm f}(\Bbbk)$ funktor serat
	netral. Maka $\omega$ berfaktor secara kanonik sebagai
	% segment-id: o014.li2.chapter9.tannakian-a.q007
	\[ \mathcal{T} \xrightarrow{\overline{\omega}} \catesubd{Comod}{\mathrm{f}}\coEnd(\omega) \xrightarrow{U} \cate{Vect}_{\mathrm{f}}(\Bbbk), \]
	dengan $\coEnd(\omega)$ memiliki struktur aljabar Hopf komutatif
	kanonik, dan $\overline\omega$ merupakan ekuivalensi kategori monoidal
	simetris.
\end{theorem}
\begin{proof}
	Karena $\mathcal{T}$ telah diketahui hingga secara lokal, pernyataan langsung
	mengikuti Teorema~\ref{prop:Tannaka-aux1} dan~\ref{prop:Tannaka-aux2}.
\end{proof}

% segment-id: o014.li2.chapter9.tannakian-a.eg001
\begin{example}
	Misalkan $H$ aljabar Hopf-$\Bbbk$.
	Proposisi~\ref{prop:Hopf-module-dual} menunjukkan bahwa kategori
	komodul berdimensi hingga $\catesubd{Comod}{\mathrm f}H$ secara alami
	menjadi kategori Tannakian netral, dengan funktor lupa sebagai funktor
	serat netral standar:
	% segment-id: o014.li2.chapter9.tannakian-a.q008
	\[ \omega: \catesubd{Comod}{\mathrm{f}}H \to \cate{Vect}_{\mathrm{f}}(\Bbbk), \quad M \mapsto (M\;\text{sebagai ruang vektor}). \]
	Teorema~\ref{prop:Tannaka-neutral} mengatakan bahwa, hingga ekuivalensi,
	keluarga contoh ini mencakup semua kategori Tannakian netral.
\end{example}

% segment-id: o014.li2.chapter9.tannakian-a.eg002
\begin{example}
	Misalkan $(\Gamma,+)$ grup komutatif dan ambil
	$\mathcal{T}=\cate{Vect}_{\Gamma,\mathrm f}(\Bbbk)$, dengan notasi dari
	Contoh~\ref{eg:Tannakian-I-graded-0} dan
	Contoh~\ref{eg:Tannakian-I-graded}. Bekali $\mathcal{T}$ dengan struktur
	kepang standar
	% segment-id: o014.li2.chapter9.tannakian-a.q009
	\begin{align*}
		c(M, N): M \otimes N & \rightiso N \otimes M \\
		x \otimes y & \mapsto y \otimes x, \quad x \in M^\gamma, \; y \in N^\eta.
	\end{align*}
	Ini menjadikan $\mathcal{T}$ kategori Tannakian netral, dengan
	$\omega_\Gamma:(M^\gamma)_{\gamma\in\Gamma}\mapsto
	\bigoplus_\gamma M^\gamma$ sebagai funktor serat netral. Pokok persoalannya ialah
	kekakuan, yang tepat merupakan versi multivariabel
	Contoh~\ref{eg:Vect-trace}. Koaljabar endomorfisme yang bersesuaian telah
	diketahui, yaitu aljabar Hopf $\Bbbk[\Gamma]$.

	Walaupun $\Gamma=\Z$ atau $\Z/2\Z$ diperbolehkan, struktur kepang di
	sini tidak melibatkan aturan tanda Koszul~\eqref{eqn:Koszul-braiding};
	jika melibatkannya, $\omega_\Gamma$ tidak akan mempertahankan struktur
	kepang. Teori kategori Tannakian tidak dapat diterapkan secara langsung
	pada kategori seperti ruang vektor super
	$\cate{Vect}^-_{\Z/2\Z,\mathrm f}(\Bbbk)$ karena kategori tersebut tidak
	memiliki funktor serat menuju $\cate{Vect}_{\mathrm f}(\Bbbk)$.
\end{example}

% segment-id: o014.li2.chapter9.tannakian-a.r001
\begin{remark}
	Beberapa pernyataan aljabar linear, yakni pernyataan tentang funktor
	serat, dapat diterjemahkan ke bahasa aljabar Hopf melalui
	Teorema~\ref{prop:Tannaka-neutral}. Sebagai contoh sederhana dan berguna,
	untuk funktor serat netral $\omega$ dan grup komutatif $(\Gamma,+)$,
	tinjau tiga jenis data berikut.
	\begin{enumerate}[(i)]
		\item Bekali setiap $\omega(X)$ dengan struktur bergradasi-$\Gamma$,
		yakni tentukan keluarga isomorfisme
		% segment-id: o014.li2.chapter9.tannakian-a.q010
		\[ \alpha_X: \omega(X) \rightiso \omega_\Gamma(FX), \quad FX \in \Obj\left( \cate{Vect}_{\Gamma, \mathrm{f}}(\Bbbk) \right), \]
		dengan $FX$ dan $\alpha_X$ funktorial dalam $X$, serta sedemikian rupa
		sehingga $\omega(\munit)\simeq\Bbbk$ dan
		$\omega(X\otimes X')\simeq\omega(X)\otimes\omega(X')$ masing-masing
		terangkat ke tingkat ruang vektor bergradasi-$\Gamma$.
		\item Tentukan funktor monoidal
		$F:\mathcal{T}\to\cate{Vect}_{\Gamma,\mathrm f}(\Bbbk)$ beserta
		isomorfisme funktor monoidal
		$\alpha:\omega\rightiso\omega_\Gamma F$.
		\item Tentukan homomorfisme bialjabar-$\Bbbk$
		$\coEnd(\omega)\to\Bbbk[\Gamma]$.
	\end{enumerate}

	Data (i) dan (ii) hanya berbeda secara redaksional. Dalam bahasa
	Definisi~\ref{def:Fib-cat}, (ii) sama dengan menentukan morfisme dalam
	$\cate{Fib}^{\otimes}(\Bbbk)$
	$(F,\alpha):(\mathcal{T},\omega)\to
	(\cate{Vect}_{\Gamma,\mathrm f}(\Bbbk),\omega_\Gamma)$. Karena itu,
	Teorema Rekonstruksi~\ref{prop:Tannaka-reconstruction} memastikan bahwa
	(ii) dan (iii) ekuivalen.

	Sebagai contoh, untuk gradasi yang paling umum $\Gamma=\Z$, data (iii)
	sama dengan memberikan homomorfisme aljabar Hopf
	$\coEnd(\omega)\to\Bbbk[T,T^{-1}]$, dengan $T$ sebuah variabel.
\end{remark}
